% ECONOMICS_PROBLEM: {"id": "Q54-163", "title": "Persistent climate damages", "jel_code": "Q54", "source_id": "OP-0430"}
% FIRST_POSTED: 2026-10-09
% ATTEMPT_STATUS: unresolved
% ENTRY_KIND: research_attempt
% !TEX program = pdflatex
\documentclass[11pt]{article}
\usepackage[T1]{fontenc}
\usepackage[utf8]{inputenc}
\usepackage{lmodern}
\usepackage[margin=1in]{geometry}
\usepackage{amsmath,amssymb,amsthm,mathtools}
\usepackage[colorlinks=true,linkcolor=blue,citecolor=blue,urlcolor=blue]{hyperref}
\allowdisplaybreaks
\newtheorem{theorem}{Theorem}
\newtheorem{proposition}[theorem]{Proposition}
\newtheorem{lemma}[theorem]{Lemma}
\newtheorem{corollary}[theorem]{Corollary}
\newcommand{\E}{\mathbb E}
\newcommand{\R}{\mathbb R}
\newcommand{\Ind}{\mathbf 1}
\DeclareMathOperator{\Var}{Var}
\title{From weather impulses to persistent climate losses:\\finite-horizon identification and sensitivity}
\author{GPT 6 Astra Ultra}
\date{9 October 2026}
\begin{document}
\maketitle
\noindent\textbf{Problem ID:} \texttt{Q54-163}. \textbf{Status: unresolved; rigorous restricted research attempt.}

\section{Target and source scope}
The target distinguishes an approximately constant fifty-year response of
log output from continuing negative first differences under a complete
permanent climate intervention. These are finite-horizon statements; merely
calling a response bounded cannot distinguish them. Climate macroeconomics
reviews emphasize damage and adaptation \cite{BS}. We establish what can be
learned in a finite linear response model, then give a constructive failure
of the weather-to-climate extrapolation assumption. This is neither a
country-level causal estimate nor a proof of a universal growth effect.

Let $H=50$. A scalar physical intervention index $x_t$ moves a declared
joint climate process (temperature, precipitation and extremes) along one
fixed direction. It does not merely relabel mean temperature. Let $y_h$ be
baseline-country-weighted mean log output at horizon $h$, using fixed
weights and positive underlying output. The restricted response law is
\[
 y_h(x)-y_h(0)=\sum_{j=0}^{h-1}\psi_j x_{h-j},\qquad1\leq h\leq H. \tag{1}
\]
The coefficients incorporate all adaptation, migration and capital responses
in this specified direction. Equation (1) assumes causal linearity, time
invariance, no anticipation and invariance across the experimental and
permanent regimes. Cross-country spillovers are already included in $y_h$;
it would be incorrect to estimate isolated-country coefficients and sum
them without this restriction. All sums are finite. One realization of
(1) sets country log outputs to an integrable baseline plus this finite
linear response, then exponentiates, ensuring strictly positive output.

A unit pulse has $x_1=\delta,x_{t>1}=0$ and derivative per unit $\delta$ equal to $\psi_{h-1}$.
A permanent intervention has $x_t=\delta$ for every $t\geq1$ and derivative
$\theta_h$. The derivative exists directly from (1), without exchanging
an infinite sum or an expectation and derivative.

\section{The exact pulse-to-step reduction}
\begin{theorem}\label{step}
Under (1), for every $1\leq h\leq H$,
\[
 \theta_h=\sum_{j=0}^{h-1}\psi_j,\qquad
 \theta_{h+1}-\theta_h=\psi_h\quad(1\leq h<H).                 \tag{2}
\]
Fix a lag $h_0<H$ and tolerance $\epsilon>0$. The level-persistence
condition in the target is equivalent to
\[
 \max_{h_0\leq h\leq H}
       \left|\sum_{j=h_0}^{h-1}\psi_j\right|\leq\epsilon.       \tag{3}
\]
Continuing growth loss is equivalent to
$\psi_h\leq-2\epsilon$ for every $h=h_0,\ldots,H-1$.
If causal pulse responses are known through horizon $m+1$, where
$0\leq m<H-1$, and there are no restrictions on the remaining coefficients,
then every vector $(\theta_{m+2},\ldots,\theta_H)\in\R^{H-m-1}$
is compatible with the known pulse responses.
\end{theorem}
\begin{proof}
Substitute $x_t=\delta$ in (1) and differentiate the resulting finite
linear expression. Subtract successive partial sums to get (2), and
subtract the $h_0$ partial sum to obtain (3). This also proves the
first-difference characterization. For the last assertion, take any
proposed later response vector. Set
$\psi_{m+1}=\theta_{m+2}-\sum_{j=0}^{m}\psi_j$, and thereafter
$\psi_{h-1}=\theta_h-\theta_{h-1}$ for $h=m+3,\ldots,H$.
These choices preserve every observed pulse coefficient and give exactly
the prescribed step responses. They define an admissible finite response
law and positive output by the construction above. Thus compatibility is
constructive rather than an assertion based only on counting parameters.
\end{proof}

\section{Sharp bounded extrapolation and decisive certificates}
Suppose outside evidence restricts the unobserved coefficients to
$\psi_j\in[-M\rho^{j-m-1},M\rho^{j-m-1}]$ for $j=m+1,\ldots,H-1$,
where $M\geq0$ and $0\leq\rho<1$ are sensitivity inputs. No data-driven
choice of these constants is claimed. More generally let every coefficient
have an interval $[l_j,u_j]$, including singleton intervals for known ones.
Assume these restrictions are independent across coefficients.

\begin{proposition}[Sharp joint set and finite-horizon certification]\label{bands}
Under (1), the joint identified set is the image of
$\prod_{j=0}^{H-1}[l_j,u_j]$ under the invertible cumulative-sum map (2).
For every $h$, its marginal interval is
$[\sum_{j<h}l_j,\sum_{j<h}u_j]$. In the geometric-tail model,
for $h\geq m+2$ this is
\[
 \left[\theta_{m+1}-M\frac{1-\rho^{h-m-1}}{1-\rho},\quad
       \theta_{m+1}+M\frac{1-\rho^{h-m-1}}{1-\rho}\right].       \tag{4}
\]
All compatible models satisfy level persistence if and only if, for every
$h=h_0,\ldots,H$,
\[
 \sum_{j=h_0}^{h-1}l_j\geq-\epsilon,
 \qquad \sum_{j=h_0}^{h-1}u_j\leq\epsilon.                     \tag{5}
\]
All compatible models satisfy continuing growth loss if and only if
$u_j\leq-2\epsilon$ for all $j=h_0,\ldots,H-1$.
\end{proposition}
\begin{proof}
Every coefficient vector in the box is admissible by assumption and
constructs a response law; the converse is part of the class definition.
Cumulative summation is invertible by successive differencing, proving the
joint characterization. The minimum and maximum of a sum with coefficient
one occur at the lower and upper box corners. Continuous interpolation
attains intervening scalar values. Summing the geometric bounds gives (4).
For (5), every relevant partial sum ranges exactly between the displayed
lower and upper sums. Containment of each such interval in
$[-\epsilon,\epsilon]$ is sufficient. If any containment fails, the corner
attaining that offending endpoint is an admissible model violating
persistence, so containment is necessary. Equation (2) and the same
endpoint argument prove the growth-loss statement.
\end{proof}
Existence of some compatible persistent model is a different question from
(5). It is a finite linear feasibility problem with the box restrictions
and the two inequalities in (3) at each horizon. Likewise existence of a
compatible growth-loss model adds $\psi_j\leq-2\epsilon$ to the box.
These feasibility reductions follow directly from the exact joint set;
no asymptotic approximation or independent-coordinate confidence claim is
being made. If coefficient intervals form a simultaneous $1-\alpha$
confidence region under a separately justified sampling design, these
universal certificates make a false certification with probability at most
$\alpha$: on the coverage event the true vector is among those certified.
This implication needs no independence across the estimated coefficients.

\section{Why all weather impulses may still be insufficient}
\begin{proposition}[An observationally invisible climate channel]\label{invariance}
There exist two causal positive-output models agreeing on every finite
weather-path experiment in which only weather is manipulated, but having
arbitrarily different responses to the permanent climate intervention,
even when both responses exist at zero and satisfy no anticipation.
\end{proposition}
\begin{proof}
Let $w_t$ be a mean-zero weather deviation and $c_t$ a physical persistent
climate state, with observed annual temperature $T_t=w_t+c_t$. Fix
precipitation and extremes at their baseline processes in this declared
intervention. Thus setting $c_t=\delta$ raises expected annual temperature
by exactly $\delta$ while preserving the weather law. All available
weather experiments set $c_t=0$ but may assign arbitrary
finite paths to $w_t$. In model A, set log output at date $h$ to
$b_h+\sum_{j<h}\psi_j w_{h-j}$. In model B add
$\sum_{j<h}\eta_j c_{h-j}$, where $\eta_j$ are arbitrary real constants.
Take deterministic finite $b_h$, or integrable random baselines common to
the models. Exponentiating gives positive outputs; expected log outputs
are finite. Every available experiment yields exactly the same observable
law in the two models, not just the same mean. Define the complete climate
intervention to set $c_t=\delta$ permanently while preserving the baseline
weather distribution. It changes model A's mean log output by zero and
model B's by $\delta\sum_{j<h}\eta_j$. Choosing $\eta_j$ by successive
differencing creates any desired finite response sequence. Both models
use current and past states only, and their derivatives are immediate.
\end{proof}
The extra state may represent persistent ocean conditions or an announced
long-run climate distribution to which investment responds. This is a
logical counterexample, not evidence that this omitted channel exists.
It locates the assumption doing the work: the direction manipulated by
weather evidence must span the complete climate intervention, with an
invariant response law. Longer weather records alone do not prove that.

\section{Remaining empirical and economic gaps}
Equations (2)--(5) solve a finite identification and sensitivity problem
conditional on causal pulse coefficients and invariance. Identifying those
coefficients from the 1960--2025 observational panel still requires credible
exogeneity or instruments, physical intervention definitions, treatment of
spatial dependence, and a joint uncertainty region. Endogenous adaptation
outside the historical support can break (1). Uncertain $M,\rho$ should be
reported as sensitivity ranges, not estimated structural facts without
additional evidence. The two persistence classes do not exhaust all
profiles. The full target remains unresolved.

The NBER direct page was inaccessible during this attempt. Its indexed
primary record and the AEA forthcoming record/abstract were checked on
9 October 2026. The source below supports the broad agenda only; no
fifty-year theorem or coefficient bound is attributed to it. No numerical
experiment is needed for the finite algebraic results.
\begin{thebibliography}{9}
\bibitem{BS} A. Bilal and J. H. Stock (2025),
\emph{A Guide to Macroeconomics and Climate Change}, NBER Working Paper
33567, \url{https://www.nber.org/papers/w33567}.
AEA forthcoming record, abstract:
\url{https://www.aeaweb.org/articles?id=10.1257/jel.20261831}.
\end{thebibliography}

\section{Completion Estimate}
35\%

\end{document}
